%\documentclass{Math_Note}
%\setlength{\parindent}{2em}

%\title{\textbf{Mathmatical Analysis Zorich Comprehensions Chapter 4}}
%\author{Buce-Ithon}
%\newdateformat{mydate}{\twodigit{\THEDAY}{ }\shortmonthname[\THEMONTH], \THEYEAR}
%\date{\today}

%\begin{document}
% Title page
%\maketitle

% Content
%\newpage
%\tableofcontents
%\newpage

% \setcounter{section}{-1}
\newpage
\section{Chapter 4.Continuous Functions}
\subsection{Comprehensions}
This chapter discusses the continuity of functions (including discontinuities when not continuous). First, we define the concepts of continuity at a point, continuity on an interval, and one-sided continuity, 
and introduces the concept of points of discontinuity ($2$ types: depending on whether the one-sided limits exist). Further, it defines the stronger concept of uniform continuity. Then, it presents the arithmetic 
operations and composition of continuous functions, and provides $3$ global properties of continuous functions (on closed intervals: 1. Bolzano-Cauchy intermediate value theorem, 2. Weierstrass extreme value theorem 
(boundedness), 3. Cantor-Heine uniform continuity theorem). Additionally, Zorich also covers the continuity of monotonic functions in the book (which is special as before) (types of discontinuities for monotonic 
functions (only the first kind) and their number (at most countable), and the relationship between continuity and the range (intermediate value property) for monotonic functions derived from this). Finally, it 
introduces the inverse function theorem, which is very intuitively understood (on closed intervals: existence $\Leftrightarrow$ injective $\Leftrightarrow$ strictly monotonic, preserving continuity).
\newline\newline
\fbox{\begin{minipage}{0.95\textwidth}
\textbf{Continuity of 2 Special fuctions: }
\begin{enumerate}
    \item[1.] Dirichlet function $\mathcal{D}(x) = \begin{cases} 1, & x \in \mathbb{Q} \\ 0, & x \in \mathbb{R} \setminus \mathbb{Q} \end{cases}$ \\
    Equivalent expression: $\forall x \in \mathbb{R}, \mathcal{D}(x) = \lim_{k \to \infty} (\lim_{j \to \infty} (\cos(k!x))^{2j})$ \\
    Properties: Every point is a discontinuity point of the second kind (no limit anywhere, discontinuous, non-differentiable); has every positive rational number as its period (thus no minimal positive period); 
    not Riemann integrable on any bounded interval, but Lebesgue integrable on $[0,1]$.
    \item[2.] Riemann function $\mathcal{R}(x) = \begin{cases} \frac{1}{n}, & x = \frac{m}{n} \in \mathbb{Q} \setminus \{0\} \\ 0, & x \in \mathbb{R} \setminus \mathbb{Q} \cup \{0\} \end{cases}$ \\
    The number-theoretic essence of the first kind discontinuity points of the Riemann function: $\forall a \in \mathbb{R}, N \in \mathbb{N}$, in $U(a)$ there are only finitely many rational numbers $m/n$ 
    (assumed to be in irreducible form) with $n < N$, i.e., $\frac{1}{n} > \frac{1}{N}$. \\
    Properties: $\mathcal{R}(x)$ is continuous at all irrational points and has first kind discontinuities at all rational points; the Riemann function is Riemann integrable.
\end{enumerate}
\end{minipage}}

\subsection{Problems Sets}
\begin{prb}
    (C4S2Q1)
\end{prb}
\fbox{\begin{minipage}{0.95\textwidth}
\begin{qsol*}
\begin{enumerate}
    \item [a)] $f \in C(A), B \subset A \implies f|_B \in C(B)$.
    \item [b)] $f : E_1 \cup E_2 \to \mathbb{R}$, satisfying $f|_{E_i} \in C(E_i), i = 1, 2 \nRightarrow f \in C(E_1 \cup E_2)$, counterexample: $f(x) = \tan x$.
    \item [c)] $\forall N > 0, \{ \frac{m}{n} : n < N, (m, n) = 1, m, n \in \mathbb{Z}^+ \}$ is a finite set, i.e. $\exists U(x; \delta) (\forall x \in (0, 1))$ such that $\forall x_1, x_2 
    \in U(x; \delta), |R(x_1) - R(x_2)| \leq \frac{1}{N} < \varepsilon$ (Cauchy convergence criterion), then $\forall x_0 \in [0, 1], \lim_{x \to x_0} R(x) = 0$. (The proof method is consistent with 
    the discussion of discontinuity points of the Riemann function in Example 2 of C4 Summary.)
\end{enumerate}
\end{qsol*}
\end{minipage}}

%\end{document}